\documentclass{article}
% Competition/workshop paper, non-anonymous (single-blind) camera-ready: the official style
% requires a track option together with [final]; [final] alone leaves the track name undefined.
\usepackage[sglblindworkshop,final]{neurips_2026}
\workshoptitle{Smart Buildings Challenge (SB26)}
\usepackage[utf8]{inputenc}
\usepackage[T1]{fontenc}
\usepackage{url}
\usepackage{booktabs}
\usepackage{amsfonts}
\usepackage{amsmath}
\usepackage{nicefrac}
\usepackage{microtype}
\usepackage{xcolor}
\usepackage{graphicx}
\usepackage{algorithm}
\usepackage{algpseudocode}
\usepackage{multirow}
\usepackage{hyperref}  % after float/algorithm so caption anchors are not duplicated
\input{results_macros}
% Figure include that degrades to a visible placeholder when the file is not built yet.
\newcommand{\figfile}[2]{\IfFileExists{#2}{\includegraphics[width=#1]{#2}}{\fbox{\parbox{#1}{\centering\textbf{??} missing figure \texttt{\detokenize{#2}}}}}}

\title{Schedule-Aware Demand-Based Reset: A Structured Rule-Based Controller for Cross-Building HVAC Transfer in the Smart Buildings Challenge}

\author{%
  {[Author Name]}\\
  {[Affiliation]}\\
  \texttt{neelsirivara@outlook.com}
}

\begin{document}

\maketitle

\begin{abstract}
The NeurIPS 2026 Smart Buildings Challenge asks for one HVAC control policy that operates
procedurally generated commercial buildings it has never seen, under 30-day weather strips
from five climates, trading occupant comfort against energy cost and carbon. We present
SADR (Schedule-Aware Demand-based Reset), a pure-numpy rule-based controller that reads the
plant topology from the building metadata and is otherwise built from facts shared by every
generated building: comfort is charged only while zones are occupied, fan power grows as
static pressure to the three-halves, the chiller has a constant coefficient of performance,
and all heat reaches a zone through a reheat coil whose discharge is capped at 308.15\,K.
SADR gates its mode on the occupancy count and schedule flags (with a self-tuning optimum
start), regulates each zone with a PI loop on \emph{net heat} translated into damper and
valve commands through the supply-air and discharge physics, and resets static pressure,
loop differential pressure and hot-water temperature from zone demand. On the 30 public
scenarios the final version scores \RFConfPubCand{} on the leaderboard metric (mean absolute
per-step reward, lower is better) against \MetricPidWater{} for the strongest organizer
baseline (\texttt{pid\_water}) under identical conditions, a paired-bootstrap ratio of
\RFPwPubRatio{} (95\% CI \RFPwPubRatioLo--\RFPwPubRatioHi; \RFPwPubWins{} wins); on \NHeldAll{}
scenarios from \NHeldAllBuild{} generated buildings and unseen strips the ratio is
\RFPwOldRatio{} (CI \RFPwOldRatioLo--\RFPwOldRatioHi). Loss decompositions drove three
revisions: coldest supply air and a deeper setback (\CandWinnerRelRed{} on the public set);
zone loops scaled with floor area and box capacity for large zones (\HeldAllVFiveRelRed{} on
the held-out buildings, CI \HeldAllVFiveRelRedLo{} to \HeldAllVFiveRelRedHi); and, after
tracing the generated buildings' larger chiller term to the per-handler energy normalization
rather than to control, a pre-registered screen of discharge headroom, pressure-aware
requests, night air and set-back ventilation whose composite lowers the loss by
\RFConfPoolRed{} (CI \RFConfPoolRedLo{} to \RFConfPoolRedHi) on \RFConfPoolN{} held-out
scenarios no decision touched, by \RFConfFreshShipRed{} on fresh buildings in the shipped
buildings' size range and by \RFConfPubRed{} on the public set. We release the controller and
report where it still loses.
\end{abstract}

\section{Introduction}

Learned building controllers have been studied for two decades, yet few run in real
buildings; the recurring obstacles are data-hungry training, fragile robustness and poor
generalization across buildings~\citep{wang2020rl}. The Smart Buildings Challenge (SB26)
isolates the last of these: a submission is a single policy class, instantiated on held-out
buildings drawn from the same procedural generator as six public ones and scored on 30-day
episodes under weather strips it has not seen~\citep{sb26kit}. The observation and action
vectors change shape with the building ($K$ air handlers and $V$ zones expose $K{+}1$ on/off
commands and $3K{+}2{+}2V$ continuous setpoints), and the simulator is Google's open-source
Smart Buildings simulator~\citep{goldfeder2023lightweight,goldfeder2025suite}.

We take the position that the right first entry in such a benchmark is a
\emph{structure-aware rule-based controller}: one that reads the plant topology from the
metadata and otherwise encodes physics every generated building shares. It transfers by
construction and shows which part of the score is reachable without learning.

\paragraph{Contributions.}
(i) An analysis of the SB26 reward and plant that yields the design facts above
(Section~\ref{sec:env}). (ii) SADR, a deterministic numpy controller (\PolicyMsOurs{}\,ms per
step) built from them (Section~\ref{sec:method}). (iii) A protocol of screens, ablation and
full-length loss decompositions that motivated three revisions, the last pre-registered and
confirmed on \RFFreshNBuild{} freshly generated buildings (Section~\ref{sec:results}). (iv) An
account of where the loss sits, why the controller lost on large zones, why the generated
buildings cost more, and where it still loses (Section~\ref{sec:analysis}), every number
generated from the result files.

\section{Environment and metric}
\label{sec:env}

\paragraph{Plant.}
Each building is a multi-storey floor plan partitioned into zones; air handlers (AHUs) each
serve up to ten zones through one VAV box per zone. An AHU mixes return air with outside
air (recirculation $r\le 0.7$) and cools the mixture to its supply-air temperature (SAT)
setpoint, so the supply is $\min(\text{mixed}, \text{SAT})$ and the AHU cannot heat. Zone
flow is $d\,F_{\max}\sqrt{P/P_{\text{design}}}$ for damper position $d\in[0.1,1]$ and static
pressure $P\in[200,2000]$\,Pa (both floored by the simulator while the unit runs). Fan power
is $\text{flow}\times P/\eta$, so the cheapest pressure for a given flow is the floor and
power grows as $P^{3/2}$ once dampers are open; the chiller's coefficient of performance is a
constant 3.5. Heat comes only from each VAV's reheat coil, fed by one gas boiler loop whose
supply temperature (310--350\,K) and differential pressure (0--20\,bar) are setpoints; the
coil's effectiveness $\min(1, C_w/C_a, (308.15-T_{\text{sa}})/(T_w - T_{\text{sa}}))$ caps the
discharge at 308.15\,K; gas is billed as heat delivered plus a tank loss.

\paragraph{Reward and metric.}
Each 5-minute step pays $0.6\,\text{(comfort shortfall)} + 0.2\,\text{(normalized energy
cost)} + 0.2\,\text{(normalized carbon)}$. The comfort shortfall is the occupancy-weighted
productivity loss (a deviation of $\Delta$\,K from the band costs $\Delta/5$ up to 5\,K and
$1+((\Delta-5)/5)^2$ beyond, normalized by total occupancy), so it is \emph{zero whenever
nobody is in the building}; occupants are present only on weekday comfort hours (arrivals
06--12, departures 13--19, none on US holidays) and only the building-wide count is
observed; the band is 294--297\,K in comfort hours and 289--298\,K otherwise. Energy is
normalized by a plant maximum that scales with $K$; in metric units one kW of electricity
costs about $2.1\times$ one kW of gas and one kW of chiller cooling about $0.6\times$. The
leaderboard metric is the mean over scenarios of
$|\text{mean per-step reward}|$, lower is better; it is dominated by the extreme climates
(the organizers' \texttt{pid\_water} loses \ExtremeRatioFairbanks$\times$ its mild-climate
mean in Fairbanks).

\paragraph{Baselines.}
The kit ships \texttt{bang\_bang}, \texttt{fixed\_setpoint}, \texttt{pid} (per-AHU PID on the
mean served-zone temperature, plant idle outside work hours) and \texttt{pid\_water}
(\texttt{pid} plus worst-zone tracking, a proportional economizer, proportional reheat and a
hot-water reset), the strongest and our main comparison; there is no learned baseline.

\section{Method: SADR}
\label{sec:method}

SADR is a single \texttt{TransferPolicy} class whose only state is per-zone integrals,
per-AHU dwell counters and the optimum-start statistics; all gains live in one parameter
dictionary (Table~\ref{tab:params}, Appendix~\ref{app:tuning}); Algorithm~\ref{alg:sadr} (Appendix~\ref{app:figs}) summarizes one step.

\paragraph{Topology from metadata.}
Index tables come from the observation field names, the positional AHU action lists, the
per-VAV action indices and the AHU-to-zone assignments; ranges and statistics from the
normalization metadata. Every action slot is written and clipped on every step; nothing depends
on the number of zones or handlers, the floor plan or the strip; non-finite inputs fall back to
schedule values; no files, no network, no randomness.

\paragraph{Modes.}
\textsc{Occupied} when the \texttt{comfort\_mode\_now} flag is set and the building is not
detected as closed; \textsc{Precondition} when unoccupied but comfort hours start within a
lead time (or \texttt{comfort\_mode\_soon}); \textsc{Setback} otherwise. ``Somebody is here''
is detected relative to the running minimum of the affinely normalized occupancy count; a
weekday on which nobody has arrived by \ParHolidayCheckH{}:00 is a holiday in the kit's
occupancy model and is treated as unoccupied until the count rises. The optimum-start lead is $\ell = \ell_0 +
\text{deficit}/(\hat{r}/s)$: the worst zone's distance from the comfort band over a
recovery-rate estimate $\hat{r}$ (K/h) re-estimated from each morning's recovery
(exponential average, weight \ParRateEma) with safety factor $s=\ParLeadSafety$, clipped to
$[\ParLeadMinH, \ParLeadMaxH]$\,h.

\paragraph{Zone loops on net heat.}
Each zone has heating and cooling targets, the sensor band shrunk by \ParMarginOcc\,K
(occupied) or \ParMarginPre\,K (preconditioning); two PI loops compute requested net heat
$q_h = k_q (T_h^\star - T)^+ + I_h$ and cooling $q_c = k_q (T - T_c^\star)^+ + I_c$ in watts
($k_q = \ParKqWPerK$\,W/K, integral gain \ParKiWPerKStep\,W/K per step); the integrals unwind
past the target, so a zone inside the band free-floats, and the larger request wins. Net
heat rather than flow because the zone-air node responds within minutes: a flow-proportional
controller with a 5-minute step oscillates, whereas a net-heat gain of the order of the air
node's conductance to the slab is stable and the integral carries the steady load.
\textsc{Setback} uses \emph{bursts} instead: a running handler pays the ventilation load of
every zone it serves, so the floor (eco band widened by \ParNightLowExt{}/\ParNightHighExt\,K)
is held by requesting \ParQSbW\,W from any zone that crosses it until it is \ParSbHystK\,K
inside, and the handler is off otherwise (holding the floor with the PI loops instead kept a
handler running for \VthreeAhuOnUnocc{} of the unoccupied steps on the coldest scenario).

\paragraph{From heat to actuators.}
A cooling zone needs capacity rate $F = q_c/(T - T_{\text{sa}})$, $T_{\text{sa}}$ being the
AHU's estimated supply temperature; a heating zone is bounded by the discharge limit, $F =
q_h/(308.15 - \delta - T)$. \emph{SAT:} a cooling handler commands the coldest value,
\ParSatCoolHi\,K. The coil load is $F\,(T_{\text{mixed}} - T_{\text{sa}}) = q_c + F\,
(T_{\text{mixed}} - T)$ and the COP is constant, so for a given $q_c$ both the ventilation term
and the fan power fall with flow and the coldest supply is cheapest (v4 reset SAT upward in
the trim-and-respond spirit of Guideline 36~\citep{ashrae2021g36,taylor2018vav}: right in a
real building, wrong here). A handler
that also heats is floored at \ParSatMixFloor\,K so its cooling zones are not overcooled into
reheat; a heating-only handler commands 305\,K so the coil idles. \emph{Static pressure} is
trim-and-respond in closed form, $P = P_{\text{design}}\max_z f_z^2$ clipped to
$[200, 2000]$\,Pa with $f_z$ each zone's requested fraction of design flow, so the most
demanding damper is fully open; dampers are $f_z/\sqrt{P/P_{\text{design}}}$. \emph{Reheat:} the discharge temperature that delivers $q_h$
at the delivered flow gives the water capacity rate $W_z = F_z\,(T_{\text{disch}} -
T_{\text{sa}})/(T_w - T_{\text{sa}})$; the loop differential pressure is reset so the largest
$W_z$ sits at \ParValveTarget{} of valve travel and each valve is $W_z/(k_v \sqrt{\Delta p})$.
An AHU runs iff a served zone requests more than \texttt{q\_on} (dwell \ParMinOnSteps{}
steps); the pump runs iff a valve is open. Recirculation is at the 0.7 cap except when a
cooling-only handler sees outside air colder than return air by more than the exhaust-fan
break-even ($0.00322$\,K/Pa$\times P$), when the outside-air fraction is solved so the mixed
air lands on the SAT target. The hot-water setpoint is 350\,K whenever heating was requested
in the last \ParHeatSeasonSteps{} steps and 310\,K otherwise: gas is billed as heat delivered,
so hotter water is strictly more useful at a tank dissipation of \TankLossExtraKw\,kW.

\paragraph{Scaling with the building (v6).}
Every VAV box in the generator has the same design flow (2\,m$^3$/s) and hot-water $k_v$, so
a 450\,m$^2$ zone has a fifth of the heating authority per square metre of a 100\,m$^2$ zone
and several times its thermal mass; v5's fixed 12\,kW cap, gains and building-wide lead were
tuned on the public zones (median \PubZoneAreaMedian\,m$^2$). v6 (i) reads each zone's floor
area from the metadata zone map (4-connected free cells in raster order reproduce the
simulator's room numbering) and scales $k_q$, $k_i$, the request cap and the setback burst with
area/\ParAreaRef\,m$^2$, clipped to $[\ParAreaScaleMin, \ParAreaScaleMax]$; (ii) bounds each
request by what the box delivers at a flow fraction $f_{\text{cap}}$ (\VSixFCapPre{} while
preconditioning or set back, \ParFCapOcc{} occupied) at the discharge limit, at 285\,K supply,
and on the water side (valve fully open at 20\,bar: asking for more raises the airflow, cools
the discharge and delivers \emph{less}); (iii) runs the optimum start per zone; (iv) raises the
setback floor of a zone whose learned rate could not recover it within the maximum lead. With
these off the code reproduces v5 bit for bit.

\paragraph{Round 4 (v7).}
\label{sec:v7}
The round-4 composite turns on five switches of the v6 code (all off reproduces v6 bit for bit),
each from a diagnosis in Section~\ref{sec:analysis}: (i) discharge margin $\delta$ from \VSixDischMargin{} to
\ParRFDischMargin\,K: the simulator caps the discharge itself, and at the cap the gas per joule of
net heat is $1 + 0.3\,(T - T_{\text{out}})/(308.15 - \delta - T)$, so headroom means less outside
air heated; (ii) an occupied zone may ask for flow above what its box delivers at the 200\,Pa
floor only in proportion to how far it is past its target (from \ParRFOccEzero\,K, full flow
\ParRFOccSpan\,K later): pressure above the floor costs fan power and enlarges every other zone's
forced minimum flow, while comfort is charged per occupant and a zone holds one or two whatever
its area; (iii) before a workday and outside heating season, set-back zones warmer than the
preconditioning target are flushed with outside air (floor pressure, coil off) while it is at
least \ParRFFlushDt\,K colder than the zone; (iv) in heating season the forced minimum flow of set-back
zones on a running handler is reheated to the zone temperature, at 1\,J of gas per joule instead of
the ventilated recovery's 1.9 later; (v) the preconditioning cap sits at the floor-pressure flow.

\paragraph{Tuning protocol.}
Parameters were first tuned by seeded random search (\NTuneConfigs{} configurations) on a
screening set of six scenarios from four public buildings (1, 2, 5, 8; one strip per station
plus a second Fairbanks strip) at 1,440 steps: the defaults scored \TuneDefaultScreen, the
best \TuneBestScreen{} and \texttt{pid\_water} \TunePidWaterScreen{} on the same runs, and a
one-at-a-time ablation (Appendix~\ref{app:tuning}) shows that doubling the integral gain to
\ParKiWPerKStep\,W/K per step is the whole gain (\AblRelDefAddKi{} on its own). Buildings 4
and 7 were never screened. That configuration (v4) was run on all 30 public scenarios at
8,640 steps, its loss decomposed (Section~\ref{sec:analysis}), and \NCandV{}
physics-motivated candidates (cold SAT; night setback widened to 8--12\,K; combinations)
screened and run at full length on the public set (Appendix~\ref{app:tuning}); the winner
became v5, and the \NHeld{} generated-building scenarios were run once after that choice. Round 3
generated \NHeldRBuild{} further buildings with the kit's generator at its default ranges
(\HeldRZoneAreaMin--\HeldRZoneAreaMax\,m$^2$ zones, \HeldRZonesMin--\HeldRZonesMax{} zones;
\NHeldR{} scenarios), screened \NScreenCands{} v6 configurations at \ScreenSteps{} steps on six
public and six held-out scenarios (Table~\ref{tab:candidates_r3}), combined the changes that
helped and shipped the composite with the best held-out screen after a full-length check under a
rule fixed beforehand (not worse than v5 on the public set, better on the held-out sets); the
\NHeldUnscr{} held-out scenarios never screened are the clean estimate. Round 4 wrote its
protocol down before running anything (\texttt{PREREG.md}): \RFNSingles{} single changes to v6
(Section~\ref{sec:v7}; three added after a first look at the screen) screened at
\RFScreenSteps{} steps on ten public scenarios and on \RFNScreenHeld{} held-out scenarios (the six
round-3 screening scenarios and \RFScrNBuild{} new screening buildings), the helpful ones combined,
and the choice confirmed at full length on the public set, on the \NHeldUnscr{} clean old
held-out scenarios and on \RFFreshNBuild{} freshly generated buildings that no decision touched.

\section{Results}
\label{sec:results}

\paragraph{Setup.}
All runs use the kit's environment loader and \texttt{SimulationTracker}, 8,640 steps per
scenario, and the kit's \texttt{score\_leaderboard.py} formula with paired bootstrap CIs over
scenarios added. Every policy in Table~\ref{tab:main} -- SADR v6, v5, v4 and the four
organizer baselines -- was run on one machine (Apple M5 Pro, 12--14 concurrent runs,
\SecPerStepOurs\,s per simulator step for v6); local \texttt{pid\_water} returns are \GapMeanPidWater{}
relative to the shipped \texttt{baseline\_returns.json} on average (range \GapMinPidWater{} to
\GapMaxPidWater), so all paired comparisons use local runs. As a cross-machine check, v4 and
\texttt{pid\_water} were also run on a Windows laptop: per-scenario relative difference
\MachineGapMean{} (max \MachineGapMax; \MachineGapPidWater{} for \texttt{pid\_water}); the
simulator's seeds are fixed, so the numbers reproduce across machines.

\begin{table}[t]
\caption{Leaderboard metric (mean over the \NScen{} public scenarios of $|$mean per-step
reward$|$; lower is better) and per-station means. Paired bootstrap 95\% CIs (\NBoot{}
resamples) are for SADR v6 minus the row (v7 against v6: Table~\ref{tab:r4}). All rows but
part of v7's run on one machine (the laptop reproduces it bit for bit); the last row is the
organizers' shipped returns. Per-scenario numbers: Appendix~\ref{app:per_scenario}.}
\label{tab:main}
\centering\scriptsize
\setlength{\tabcolsep}{3pt}
\begin{tabular}{lcccccc c}
\toprule
policy & all & Fairbanks & Denver & Atlanta & Phoenix & Singapore & $\Delta$ vs v6 [95\% CI] \\
\midrule
SADR v7 (final, round 4) & \textbf{\RFConfPubCand} & \RFConfPubCandFairbanks & \RFConfPubCandDenver & \RFConfPubCandAtlanta & \RFConfPubCandPhoenix & \RFConfPubCandSingapore & Table~\ref{tab:r4} \\
SADR v6 (round 3) & \MetricOurs & \MetricOursFairbanks & \MetricOursDenver & \MetricOursAtlanta & \MetricOursPhoenix & \MetricOursSingapore & -- \\
SADR v5 & \MetricOursVFive & \MetricOursVFiveFairbanks & \MetricOursVFiveDenver & \MetricOursVFiveAtlanta & \MetricOursVFivePhoenix & \MetricOursVFiveSingapore & \DiffOursVFive{} [\DiffLoOursVFive, \DiffHiOursVFive] \\
SADR v4 & \MetricOursVFour & \MetricOursVFourFairbanks & \MetricOursVFourDenver & \MetricOursVFourAtlanta & \MetricOursVFourPhoenix & \MetricOursVFourSingapore & \DiffOursVFour{} [\DiffLoOursVFour, \DiffHiOursVFour] \\
pid\_water & \MetricPidWater & \MetricPidWaterFairbanks & \MetricPidWaterDenver & \MetricPidWaterAtlanta & \MetricPidWaterPhoenix & \MetricPidWaterSingapore & \DiffPidWater{} [\DiffLoPidWater, \DiffHiPidWater] \\
pid & \MetricPid & \MetricPidFairbanks & \MetricPidDenver & \MetricPidAtlanta & \MetricPidPhoenix & \MetricPidSingapore & \DiffPid{} [\DiffLoPid, \DiffHiPid] \\
fixed\_setpoint & \MetricFixedSp & \MetricFixedSpFairbanks & \MetricFixedSpDenver & \MetricFixedSpAtlanta & \MetricFixedSpPhoenix & \MetricFixedSpSingapore & \DiffFixedSp{} [\DiffLoFixedSp, \DiffHiFixedSp] \\
bang\_bang & \MetricBangBang & \MetricBangBangFairbanks & \MetricBangBangDenver & \MetricBangBangAtlanta & \MetricBangBangPhoenix & \MetricBangBangSingapore & \DiffBangBang{} [\DiffLoBangBang, \DiffHiBangBang] \\
pid\_water (organizer values) & \MetricPidWaterOff & \MetricPidWaterOffFairbanks & \MetricPidWaterOffDenver & \MetricPidWaterOffAtlanta & \MetricPidWaterOffPhoenix & \MetricPidWaterOffSingapore & \DiffPidWaterOff{} [\DiffLoPidWaterOff, \DiffHiPidWaterOff] \\
\bottomrule
\end{tabular}
\end{table}

\paragraph{Main result.}
SADR v7 reaches \RFConfPubCand{} against \MetricPidWater{} for \texttt{pid\_water}
(Table~\ref{tab:main}), a ratio of \RFPwPubRatio{} (95\% CI \RFPwPubRatioLo{} to
\RFPwPubRatioHi; wins \RFPwPubWins); v6 reached \MetricOurs{} (reduction \RelRedPidWater, CI
\RelRedLoPidWater{} to \RelRedHiPidWater; Wilcoxon $p\PWilcoxonPidWater$; per scenario in
Appendix~\ref{app:per_scenario}). On the public set v4 to v5
cut the metric by \CandWinnerRelRed{} (wins \CandWinnerWins/\NScen, mostly in the cooling
climates) and v5 to v6 by \RelRedOursVFive{} (CI \RelRedLoOursVFive{} to \RelRedHiOursVFive;
wins \WinsOursVFive/\NScen), small because the public zones are small. For scale, the public
Kaggle leaderboard (\LbDate; \LbNTeams{} teams; unseen buildings and strips, so not comparable)
stood at \LbRankOne{} for rank 1 and \LbRankEight{} for rank 8, the last prize rank.

\paragraph{Where the loss goes.}
Logging the plant energy rates every step splits the loss into comfort and the fan, chiller,
pump and boiler shares of the energy and carbon terms (Figure~\ref{fig:decomp}; by hour in
Appendix~\ref{app:figs}). On the public set v5 pays \DecPubVFiveGasShare{} of its loss for gas,
\DecPubVFiveCoolingShare{} for the chiller, \DecPubVFiveFanShare{} for fans and
\DecPubVFiveComfortShare{} for comfort (Fairbanks: gas \DecPubVFiveGasShareFairbanks);
\DecPubVFiveModePreShare{} is spent preconditioning and \DecPubVFiveModeSbShare{} set back;
\texttt{pid\_water} pays \RegretSharePidWater{} of its loss for comfort. v6 lowers comfort from
\DecPubVFiveComfort{} to \DecPubOursComfort{} and gas from \DecPubVFiveGas{} to \DecPubOursGas{}
per step.

\paragraph{Held-out buildings.}
v5 and \texttt{pid\_water} are fixed policies; v6 was run on both held-out sets once, after the
choice (Figure~\ref{fig:heldout}; every scenario in Appendix~\ref{app:per_scenario}).
Pooled over the \NHeldAll{} scenarios of the \NHeldAllBuild{} generated buildings, v6 scores
\HeldAllMetricOurs{} against \HeldAllMetricOursVFive{} for v5 (difference
\HeldAllVFiveDiff, CI [\HeldAllVFiveDiffLo, \HeldAllVFiveDiffHi]; wins
\HeldAllVFiveWins/\NHeldAll; Wilcoxon $p\HeldAllVFivePWilcoxon$) and \HeldAllMetricPidWater{}
for \texttt{pid\_water} (ratio \HeldAllPidWaterRatio, CI
\HeldAllPidWaterRatioLo--\HeldAllPidWaterRatioHi, wins \HeldAllPidWaterWins/\NHeldAll; v5:
\HeldAllRatioVFivePidWater). On the \NHeldUnscr{} scenarios never screened the reduction from
v5 is \HeldUnscrVFiveRelRed{} (CI \HeldUnscrVFiveRelRedLo{} to \HeldUnscrVFiveRelRedHi; wins
\HeldUnscrVFiveWins/\NHeldUnscr): buildings 0501/0502 go from \HeldMetricOursVFive{} to
\HeldMetricOurs{} and the \NHeldRBuild{} round-3 buildings (median zone
\HeldRZoneAreaMedian\,m$^2$) from \HeldRMetricOursVFive{} to \HeldRMetricOurs.

\paragraph{Tuning and candidates.}
The random-search spread (\TuneBestScreen--\TuneWorstScreen, Appendix~\ref{app:figs}) is small
against the gap to \texttt{pid\_water} (\TunePidWaterScreen): the gains come from the
structure. The best of the \NCandV{} v5 candidates reached \CandWinnerMetric{} at full length
(v4: \CandWinnerVFourMetric; Appendix~\ref{app:tuning}). The round-3 ledger
(Table~\ref{tab:candidates_r3}) keeps the failures: against the v6 defaults \NRThreeWorse{} of
\NRThreeVariants{} variants were worse on a screen (without the area scaling the held-out gain over
v5 disappears; the preconditioning latch lost), \NRThreeNs{} were not significant and
\NRThreeRuleMet{} met the rule; at full length the v6 defaults alone reach \FullCZeroPub{} on the
public set and the shipped variant \FullShipPub.

\begin{table}[t]
\caption{Round 4 at full length (8,640 steps): the composite v7 (\RFChosen) against the shipped
v6 on the same scenarios; $\Delta$ = paired mean difference in $10^{-5}$ per step with its
bootstrap 95\% CI (\NBoot{} resamples; negative favours v7). Old clean = the round-2/3 held-out
scenarios no screen touched; fresh = \RFFreshNBuild{} buildings generated for round 4 and run only
for this table, half in the shipped six's parameter ranges, half in the kit's default ranges.}
\label{tab:r4}
\centering\scriptsize
\setlength{\tabcolsep}{4pt}
\begin{tabular}{lrrrrrr}
\toprule
set & $n$ & v6 & v7 & $\Delta$ [95\% CI] & relative [95\% CI] & wins \\
\midrule
public & \RFConfPubN & \RFConfPubVSix & \RFConfPubCand & \RFConfPubDiff{} [\RFConfPubLo, \RFConfPubHi] & \RFConfPubRel{} [\RFConfPubRelLo, \RFConfPubRelHi] & \RFConfPubWins \\
old held-out, clean & \RFConfOldCleanN & \RFConfOldCleanVSix & \RFConfOldCleanCand & \RFConfOldCleanDiff{} [\RFConfOldCleanLo, \RFConfOldCleanHi] & \RFConfOldCleanRel{} [\RFConfOldCleanRelLo, \RFConfOldCleanRelHi] & \RFConfOldCleanWins \\
fresh, shipped ranges & \RFConfFreshShipN & \RFConfFreshShipVSix & \RFConfFreshShipCand & \RFConfFreshShipDiff{} [\RFConfFreshShipLo, \RFConfFreshShipHi] & \RFConfFreshShipRel{} [\RFConfFreshShipRelLo, \RFConfFreshShipRelHi] & \RFConfFreshShipWins \\
fresh, default ranges & \RFConfFreshDefN & \RFConfFreshDefVSix & \RFConfFreshDefCand & \RFConfFreshDefDiff{} [\RFConfFreshDefLo, \RFConfFreshDefHi] & \RFConfFreshDefRel{} [\RFConfFreshDefRelLo, \RFConfFreshDefRelHi] & \RFConfFreshDefWins \\
fresh, all & \RFConfFreshN & \RFConfFreshVSix & \RFConfFreshCand & \RFConfFreshDiff{} [\RFConfFreshLo, \RFConfFreshHi] & \RFConfFreshRel{} [\RFConfFreshRelLo, \RFConfFreshRelHi] & \RFConfFreshWins \\
clean held-out, pooled & \RFConfPoolN & \RFConfPoolVSix & \RFConfPoolCand & \RFConfPoolDiff{} [\RFConfPoolLo, \RFConfPoolHi] & \RFConfPoolRel{} [\RFConfPoolRelLo, \RFConfPoolRelHi] & \RFConfPoolWins \\
\bottomrule
\end{tabular}
\end{table}

\paragraph{Round 4.}
Of the \RFNSingles{} single changes (Table~\ref{tab:candidates_r4}, Appendix~\ref{app:tuning}),
\RFNHelps{} helped on the held-out screen without hurting the public one, \RFNPubOnly{} on the
public screen only, \RFNNs{} were not significant and \RFNWorse{} were worse: preconditioning caps
below the floor-pressure flow, which only slow the recovery, and the night flush in heating
season, whose heat had to be bought back on a winter Fairbanks strip (mean outside air
$\RFFlushFailOatC^\circ$C). The composite of the best change per family (margin 0, the milder
occupied cap, the guarded flush, set-back reheat, and the floor-pressure cap, which helped on the
public screen only) beat v6 by \RFCFourteenHeldD{} (CI
[\RFCFourteenHeldLo, \RFCFourteenHeldHi], wins \RFCFourteenHeldWins) and \RFCFourteenPubD{} (CI
[\RFCFourteenPubLo, \RFCFourteenPubHi], wins \RFCFourteenPubWins) $\times10^{-5}$ per step on the
two screens. At full length (Table~\ref{tab:r4}) v7 beats v6 on every set, by \RFConfPubRed{}
on the public set (CI \RFConfPubRedLo{} to \RFConfPubRedHi; wins \RFConfPubWins),
\RFConfOldCleanRed{} on the clean old held-out scenarios and \RFConfFreshRed{} on the fresh ones
(\RFConfFreshShipRed{} in the shipped six's ranges, \RFConfFreshDefRed{} in the default ranges),
so the pre-registered rule is \RFShipRuleMet{} and v7 is the submission. The gain is mostly gas (public:
\RFDecPubGasVSix{} to \RFDecPubGasVSeven{} per step; Fairbanks \RFDecPubFairbanksRel), then
chiller and fan on the generated buildings (old held-out: \RFDecOldCoolingDiff{} and
\RFDecOldFanDiff{} $\times10^{-5}$), while comfort rises slightly (\RFDecOldComfortVSix{} to
\RFDecOldComfortVSeven): a few per cent everywhere, not a new regime.

\section{Analysis}
\label{sec:analysis}

\paragraph{Why it wins.}
For \texttt{pid\_water} comfort regret is \RegretSharePidWater{} of the loss and
\OccOutPidWater{} of occupied zone-steps lie outside the band (mean worst-zone deviation
\OccWorstDevPidWater\,K): its loop regulates one temperature per handler while zones drift,
and its start is a step at 06:00 (Figure~\ref{fig:hours}); SADR v6 keeps \OccOutOurs{} of
occupied zone-steps out of band (comfort \RegretShareOurs{} of its loss) and in the extreme
climates uses far less energy: floor pressure with wide dampers, hot water at 350\,K with the
least flow the discharge cap allows, recirculation at the cap, and no conditioning while nobody
is charged for comfort.

\paragraph{Where the loss sits.}
For v6, \UnoccShareOurs{} of the loss is incurred on unoccupied steps
(\UnoccShareOursFairbanks{} in Fairbanks), most at hour \PeakLossHourOurs{}
(\PeakLossShareOurs{} of the day), the warm-up before 06:00. Fairbanks is gas: heating with a
mandatory 30\% outside-air fraction at $-20^\circ$C costs about $1.9$\,J of gas per joule of
net heat at the discharge cap, and no plant here heats any other way. Phoenix and Singapore are
chiller electricity, where the coldest supply at the lowest flow is cheapest (coil load $q_c +
F\,(T_{\text{mixed}} - T)$, constant COP).

\paragraph{Why v5 lost on large zones.}
On building 0502 comfort was \DecHeldVFiveComfortShareBgenb{} of v5's loss (public:
\DecPubVFiveComfortShare) and its two 350--430\,m$^2$ zones sat 2--4\,K below the band until
noon in the cold strips: the 12\,kW cap was binding, the single lead saturated at
\ParLeadMaxH\,h, and the handlers ran all night into walls that absorb the heat (ten-day test:
\TenDayVFiveScore{} for v5, \TenDayVSixBScore{} for v6). On the round-3 buildings the
largest-quartile recovery takes \DecHeldRVFiveRecoveryBig\,h for v5 and
\DecHeldROursRecoveryBig\,h for v6, and comfort is \DecHeldRVFiveComfortShare{} of v5's loss
against \DecHeldROursComfortShare{} of v6's; the price is fan energy (\DecHeldRVFiveFan{} to
\DecHeldROursFan{} per step).

\paragraph{What round 4 found.}
v6's chiller term on the generated buildings (\DecHeldROursCooling{} per step against
\DecPubOursCooling{} on the public set) is not a control defect: per square metre of zone they cool
no more than the public buildings (\RFCoolWmRThree{} against \RFCoolWmPub\,W/m$^2$ of chiller
electricity in the Phoenix and Singapore strips), but energy is normalized per handler, a handler
serves at most ten zones, and with \RFAreaPerAhuRThree{} instead of \RFAreaPerAhuPub\,m$^2$ of zone
per handler the same physics costs twice as much (\RFCoolNormRThree{} against \RFCoolNormPub{} of
the normalizer). In those strips the economizer carries \RFEconHotRThree{} of the coil load and the
load is flat through comfort hours; in the Denver and Atlanta strips that cool, night air is at
least 3\,K colder than the zones on \RFFlushMildRThree{} of unoccupied steps. What control can still reach is fan
pressure above the floor (\RFFanAboveRThree{} per step, \RFFanAboveRThreeShare{} of the metric;
\RFFanAboveMixRThreeShare{} of it on handlers that heat and cool at once) and ventilation in the
cold strips: the outside-air heat of running handlers equals \RFVentGasPub{} of the gas in the
public Fairbanks strips below $0^\circ$C, and idle zones chilled by their forced minimum flow while
set back or preconditioning another \RFIdleUnoccGasPub. Last, the kit's default generator ranges
do not reproduce the six shipped buildings (seed 0 gives \RFSeedZeroDefaultSize{} ft against
\texttt{building\_train\_1}'s \RFTrainOneSize), while the kit says the test buildings share the
shipped six's ranges: the round-3 held-out buildings are larger than the likely test buildings, so
round 4 reports both families.

\paragraph{Where it loses.}
SADR v7 is worse than \texttt{pid\_water} on \RFNScenWorseCand{} of \NScen{} public scenarios;
its largest residual loss is on \RFMaxCandScen{} (\RFMaxCandScore), almost all gas. What remains is the
morning recovery in the coldest strips, which the lead keeps cheap rather than perfect
(Figure~\ref{fig:hours}), and on the largest generated zones a box's 2\,m$^3$/s and $k_v$ that are
too small for a cold morning: capacity, not control. A learned policy's advantage should be looked
for there.

\section{Limitations}
SADR encodes constants of this simulator (discharge limit, valve coefficient, constant COP,
exhaust-fan break-even) that would be commissioning parameters in a real building, where cold
SAT would also be weighed against humidity. Candidate selection used the public set at full
length (v5) and screens on both sets (v6, v7), so the public number is optimistic; the \RFConfPoolN{} held-out
scenarios that no decision touched (\RFFreshNBuild{} buildings generated only for the final check)
are the honest ones, still a small sample of the generator, and the Kaggle test buildings may
differ further. Model predictive control or a learned residual on top of SADR is the natural
next step.

\section{Reproducibility and disclosure}
Code: \url{https://github.com/neelsirivara/neurips2026-smart-buildings} (kit and time series
excluded): \texttt{submission/transfer\_policy.py} is the policy; \texttt{experiments/src/}
generates every number and figure (\texttt{finalize.sh}; generated buildings:
\texttt{make\_sets\_r4.sh} and \texttt{STATUS.md}); simulator seeds fixed; Python 3.12.15, CPU only. An LLM coding assistant (Claude) wrote code and
drafted this text; every number is verified in \texttt{paper/AUDIT.md}. No external data or
models were used.

\bibliographystyle{plainnat}
\bibliography{refs}

\clearpage
\appendix
\section{Per-scenario results}
\label{app:per_scenario}
Tables~\ref{tab:per_scenario}, \ref{tab:per_scenario_held} and~\ref{tab:per_scenario_held_r3}
list every scenario's total return and leaderboard score (SADR v6), and
Table~\ref{tab:per_scenario_fresh} the round-4 confirmation set; all values are generated from the
result files by \texttt{make\_macros.py}.
\input{per_scenario_table}
\input{per_scenario_table_r4}
\begin{figure}[h]
\centering
\figfile{0.92\linewidth}{figures/fig_round4_buildings.pdf}
\caption{Round 4: change of the metric from v6 to v7 on every generated held-out building
(diamonds: the building's five strips pooled; dots: single scenarios; negative is better), each
run once at full length; buildings grouped by set and ordered by zone count.}
\label{fig:r4buildings}
\end{figure}
\begin{figure}[h]
\centering
\figfile{0.86\linewidth}{figures/fig_scores_by_scenario.pdf}
\caption{Per-scenario leaderboard score (log scale, lower is better), SADR v6 vs
\texttt{pid\_water} on the same machine, grouped by station and labelled by building (b =
public, g = generated held-out). Left: the 30 public scenarios; right: the \NHeld{}
held-out scenarios of buildings 0501 and 0502.}
\label{fig:scores}
\end{figure}

\section{Parameters, ablation and candidates}
\label{app:tuning}
\begin{table}[h]
\caption{Final SADR parameters (v7); values that changed from v4 to v5 are marked; the fifth row
is new in v6 (round 3), the discharge margin (v6: \VSixDischMargin\,K), the preconditioning cap
(v6: \VSixFCapPre) and the last row are round 4. The full dictionary is in the code.}
\label{tab:params}
\centering\footnotesize
\begin{tabular}{llll}
\toprule
occupied margin & \ParMarginOcc\,K & precondition margin & \ParMarginPre\,K \\
setback widening (low/high) & \ParNightLowExt{}/\ParNightHighExt\,K (v4: \VFourNightLowExt/\VFourNightHighExt) & $k_q$, $k_i$ & \ParKqWPerK, \ParKiWPerKStep\,W/K \\
cooling SAT / mixed floor & \ParSatCoolHi{}/\ParSatMixFloor\,K (v4 ceiling: \VFourSatCoolHi) & discharge margin & \ParDischMarginK\,K \\
lead safety / rate$_0$ & \ParLeadSafety{} / \ParRateZeroKph\,K/h & lead range & \ParLeadMinH--\ParLeadMaxH\,h \\
heating-season memory & \ParHeatSeasonSteps{} steps & AHU dwell & \ParMinOnSteps{} steps \\
area scaling (v6) & per \ParAreaRef\,m$^2$, clip $[\ParAreaScaleMin, \ParAreaScaleMax]$ & flow cap $f_{\text{cap}}$ (v6) & \ParFCapPre{} pre/setback, \ParFCapOcc{} occ. \\
occupied pressure cap (v7) & from \ParRFOccEzero\,K, span \ParRFOccSpan\,K & night flush (v7) & $\ge$\ParRFFlushDt\,K, $\le$\ParRFFlushWindow\,h before work \\
\bottomrule
\end{tabular}
\end{table}

Table~\ref{tab:ablation} reverts each tuned v4 value on its own (and adds each on its own to
the defaults) on the screening set; Table~\ref{tab:candidates} lists every v5 candidate at
full length on the 30 public scenarios against v4 on the same machine;
Table~\ref{tab:candidates_r3} lists every round-3 (v6) variant on both screening sets against
v5 and Table~\ref{tab:candidates_r4} every round-4 candidate against v6. The ideas tried and
rejected, with their numbers, are in \texttt{experiments/results/round3/candidates.json} and
\texttt{round4/candidates.json}.
\input{tuning_tables}

\section{Algorithm and additional figures}
\label{app:figs}
\begin{figure}[h]
\centering
\begin{minipage}[t]{0.5\linewidth}\centering
\figfile{\linewidth}{figures/fig_loss_decomp.pdf}
\caption{Per-station decomposition of the per-step loss (metric units) into comfort, fan,
chiller, pump and gas on the 30 public scenarios, SADR v5 and v6.}
\label{fig:decomp}
\end{minipage}\hfill
\begin{minipage}[t]{0.44\linewidth}\centering
\figfile{\linewidth}{figures/fig_heldout_buildings.pdf}
\caption{Metric per generated held-out building (log scale, lower is better; ordered by zone
count): \texttt{pid\_water}, SADR v5 and v6, each run once at full length on the building's
five strips.}
\label{fig:heldout}
\end{minipage}
\end{figure}
\begin{algorithm}[h]
\caption{SADR, one step (vectorized over zones $z$ and handlers $k$)}
\label{alg:sadr}
\begin{algorithmic}[1]
\State mode $\gets$ \textsc{Occupied} / \textsc{Precondition} / \textsc{Setback} from flags, time, occupancy count, lead $\ell$
\State $(T_h^\star, T_c^\star)_z \gets$ band for the mode; update $I_h, I_c$; $q_h, q_c \gets$ PI requests (W), or setback bursts
\State run$_k \gets$ hysteresis on $\max_{z\in k}(q_h, q_c)$
\State SAT$_k \gets$ \ParSatCoolHi\,K if cooling (\ParSatMixFloor\,K if also heating), 305\,K otherwise; recirc$_k \gets$ economizer
\State $F_z \gets q_c/(T_z - T_{\text{sa},k})$ or $q_h/(308.15-\delta-T_z)$; $P_k \gets P_{\text{design}}\max_z f_z^2$; damper$_z \gets f_z/\sqrt{P_k/P_{\text{design}}}$
\State $W_z \gets F_z (T_{\text{disch},z} - T_{\text{sa}})/(T_w - T_{\text{sa}})$; $\Delta p \gets (\max_z W_z / (0.9 k_v))^2$; valve$_z \gets W_z/(k_v\sqrt{\Delta p})$
\State $T_w^{\text{sp}} \gets 350$\,K if heating season else 310\,K; pump $\gets$ any valve open
\State scatter all setpoints into the normalized action vectors; clip
\end{algorithmic}
\end{algorithm}
\begin{figure}[h]
\centering
\figfile{\linewidth}{figures/fig_loss_by_hour.pdf}
\caption{Mean per-step loss by hour of day over the \NAnalysis{} public scenarios, split
into comfort regret and energy\,+\,carbon; comfort hours (06--19) shaded. Left: SADR v6;
right: \texttt{pid\_water}.}
\label{fig:hours}
\end{figure}
\begin{figure}[h]
\centering
\figfile{\linewidth}{figures/fig_trace_building_train_8_fairbanks_8_day8.pdf}
\caption{One weekday on \texttt{building\_train\_8/fairbanks\_8}: zone-temperature
envelope (min--max shaded, median line) against the comfort band (green wash), with the
building occupancy count shaded in grey. Left: SADR v6; right: \texttt{pid\_water}.}
\label{fig:trace}
\end{figure}
\begin{figure}[h]
\centering
\figfile{0.5\linewidth}{figures/fig_tuning.pdf}
\caption{Screening score (six scenarios, five days) of every random-search configuration,
sorted, with the design defaults highlighted and \texttt{pid\_water} on the same runs.}
\label{fig:tuning}
\end{figure}

\end{document}
